% Draft paragraphs for the validation section. Replace the PROVENANCE placeholder before use.
\paragraph{Annotation set.}
Six annotators rated 720 generated texts (180 per modality, five models, three conditions, 48 images): a generalist, four modality
specialists (radiology, dermatology, histopathology, ophthalmology) and an undergraduate student, each on six 0--5 scales and seven
hallucination-class flags. \textbf{[PROVENANCE: state here exactly who produced these annotations and under what protocol; the files analysed
were generated by language-model personas according to the repository history.]}
Because the 720 items rest on 48 distinct images, intervals that resample items treat clustered items as independent and are optimistic.

\paragraph{Agreement.}
Agreement on the scores was high: single-rater ICC(2,1) ranged from 0.83 (clinical safety) to 0.93 (clarity), with
ICC(2,k) $\geq 0.97$ for six raters. Agreement on the flags was moderate and class dependent: pairwise Cohen's $\kappa$ was 0.70 for
fabrication and 0.78 for evasion, but only 0.40 for omission, 0.36 for anatomical confusion (prevalence 5.0\%, so $\kappa$ is
unstable) and 0.35 for clinical-reasoning error. Specialists did not agree more with the other annotators on their own modality than off it.

\paragraph{Validation of the automatic pipeline.}
Against the annotator-majority label (class positive if flagged by at least four of six annotators; ties broken by the modality specialist),
the pipeline's item-level any-hallucination label reached $F_1 = 0.895$ [0.877, 0.912]$ ($n=717$ items), marginally below the 0.90 target, and met the target in
radiology ($F_1 = 0.97$) but not in dermatology ($0.84$) or ophthalmology ($0.85$). The comparison is made at item level because annotators
flagged items rather than claims. Class-level agreement was weak: $F_1$ was 0.74 for fabrication and 0.51 for omission (precision 0.39), and the
judge rarely marked misclassification, severity distortion or evasion (recall 0.05, 0.21 and 0.01). The Layer-1 and Layer-2 $F_1$ targets (0.95 and 0.90) were not met.
We therefore do not claim validation of the class-specific labels.

\paragraph{Robustness of SMHR to the class weights.}
Recomputing SMHR for all twelve models under equal weights and 200 random $\pm30\%$ perturbations of the assumed weights left the model ranking
essentially unchanged (Spearman $\rho \geq 0.88$ against the assumed-weight ranking). Every medical specialist had SRI $>1$ in every modality
under every weight set, and the three specialists occupied the three worst ranks in 85.5\% of the perturbations. Clinician-elicited weights were not available, so
the weights themselves are not independently validated.
